\subsection{Plataformas utilizadas y criterios de uso}\label{sec:plataformas}

Para entender el procedimiento conviene distinguir qué aporta cada herramienta. Scamper y RIPE Atlas generan las mediciones; Wireshark permite revisar los paquetes que producen; NDEx sirve para compartir los grafos construidos a partir de los resultados. Las usé en distintas etapas del mismo trabajo: obtener rutas, contrastar las respuestas y publicar una representación que otra persona pueda consultar.

\subsubsection{Scamper: medir desde la conexión de origen}

Scamper es una herramienta de medición activa: envía sondas a la red y procesa las respuestas para estudiar su topología y comportamiento. En esta tarea lo ejecuté por consola desde mi conexión domiciliaria, seleccionando el destino, el método y los parámetros de emisión. Conservé el resultado original en WARTS y una conversión a JSON para procesarlo después \cite{ref1}.

Su utilidad está en el control del experimento. Con \texttt{trace} comparé seis métodos y con \texttt{tracelb} exploré la información adicional que podía aparecer al considerar balanceo de carga. Eso permite estudiar qué cambia al modificar la forma de sondear desde un mismo origen. Su límite práctico es el punto de observación: una ejecución local describe lo que se observa desde casa; para medir desde otro extremo hace falta disponer allí de otra instancia o de una sonda remota.

\subsubsection{RIPE Atlas: medir desde sondas distribuidas}

RIPE Atlas es una plataforma de medición de conectividad que reúne sondas distribuidas en distintas redes. El usuario define una medición y selecciona las sondas que deben ejecutarla. Las solicitudes pueden gestionarse desde la web o mediante la API, y los resultados quedan asociados a un identificador que permite recuperarlos y citarlos. La ejecución depende de la disponibilidad de las sondas y de los créditos y límites aplicables \cite{plat_atlas,plat_atlas_udm}.

En el proyecto instalé una sonda de software en Docker, con ID 1018487, y solicité traceroutes UDP, ICMP y TCP por API. Esa sonda permitió medir la ida desde mi conexión; las sondas de los destinos permitieron medir la vuelta hacia mi IPv4 pública. Hay dos identificadores distintos: el ID de sonda indica quién mide y el ID de medición identifica el experimento solicitado.

Frente a Scamper, su principal ventaja para esta tarea es acceder al otro extremo sin administrar un equipo en cada destino. A cambio, hay que comprobar que la solicitud realmente produjo un resultado. Además, una captura hecha en casa no muestra todo el tráfico emitido por una sonda remota. Por eso guardé tanto las definiciones como los resultados originales, sus IDs y los tiempos de ejecución.

\subsubsection{Wireshark, Dumpcap y TShark: revisar los paquetes}

Wireshark es un analizador de paquetes. Permite abrir una captura, filtrar tráfico e inspeccionar los encabezados y bytes de cada paquete. En este estudio lo usé para observar TTL, puertos, flags, secuencias, checksums y payload, y para contrastar esas observaciones con la salida de las herramientas de medición \cite{plat_wireshark}.

Dentro del mismo conjunto de herramientas, Dumpcap realizó la captura en PCAPNG y TShark permitió extraer campos por consola para repetir el análisis de forma sistemática \cite{plat_dumpcap,ref2}. La interfaz gráfica facilitó la inspección de casos concretos; la exportación automatizada permitió comparar las rondas. Esta etapa fue especialmente útil con TCP: dejó distinguir las respuestas presentes en la captura de los saltos que Scamper incorporó a WARTS. La evidencia sigue limitada al tráfico visible en la interfaz donde se capturó.

\subsubsection{NDEx: representar y compartir lo observado}

NDEx, \emph{Network Data Exchange}, permite almacenar, visualizar y compartir redes representadas como grafos. Recibe nodos, enlaces y atributos, y ofrece una ficha con su descripción e identificador UUID. También dispone de una API para trabajar con esas redes \cite{ref5,plat_ndex}.

En esta tarea lo usé después del procesamiento: construí un grafo dirigido por destino, lo exporté a CX2 y lo publiqué en NDEx. En los nodos y enlaces quedaron representadas las interfaces observadas, los huecos y los sentidos de las mediciones. La descripción conserva los dos IDs RIPE utilizados para que el lector pueda volver a los resultados originales.

Su función es distinta de la de RIPE Atlas. Atlas aporta las mediciones de conectividad; NDEx permite presentar y compartir el grafo derivado de ellas. El UUID identifica ese grafo publicado, no una medición ni una sonda. Publicarlo facilita la revisión, pero la validez de las adyacencias sigue dependiendo de los datos y de las decisiones con que se construyó.

\subsubsection{Herramientas de apoyo: contexto y entorno de ejecución}

RIPEstat y BGP HE aportaron contexto para interpretar las IP: prefijos anunciados y asociaciones con sistemas autónomos. Esa información complementa las rutas observadas y no reemplaza las respuestas del traceroute \cite{ref4,ref7}.

Docker permitió ejecutar las herramientas y la sonda en contenedores. Docker Compose reunió la configuración de servicios, red y volúmenes en un archivo YAML, facilitando iniciarlos con los mismos parámetros \cite{plat_compose}. En la sonda se utilizó la red del host para evitar añadir una red de puente de Docker al recorrido. Los contenedores ayudan a conservar el entorno; la reproducibilidad también requiere guardar parámetros, horarios y resultados.

\subsubsection{Comparación práctica de uso}

La siguiente comparación resume cómo se usaron las herramientas en este proyecto. La elección depende de la pregunta: medir una ruta, observar sus paquetes, interpretarla o compartirla.

\begingroup
\fontsize{9}{11}\selectfont
\addtocounter{table}{-1}
\begin{longtable}{p{2.1cm}p{3.4cm}p{5.0cm}p{4.2cm}}
\toprule
\textbf{Herramienta} & \textbf{Forma de uso} & \textbf{Aporte al trabajo} & \textbf{Condición o límite} \\ \midrule
\endfirsthead
\toprule
\textbf{Herramienta} & \textbf{Forma de uso} & \textbf{Aporte al trabajo} & \textbf{Condición o límite} \\ \midrule
\endhead
Scamper & Consola y scripts desde casa. & Elegir métodos y conservar rutas en WARTS/JSON. & La medición parte del equipo donde se ejecuta. \\
RIPE Atlas & API para solicitar y recuperar mediciones; web para consultar sondas y resultados. & Medir desde la sonda local y desde las sondas remotas. & Disponibilidad, créditos y participación efectiva. \\
Wireshark y herramientas asociadas & Inspección gráfica, captura y exportación por consola. & Revisar encabezados y payload; contrastar paquetes con resultados. & Solo se observa el tráfico accesible al punto de captura. \\
NDEx & Importación del CX2 y consulta del grafo en la web. & Publicar tres grafos con UUID y referencias RIPE. & La interpretación depende del grafo que se carga. \\
RIPEstat y BGP HE & Consulta de información de IP, prefijo y ASN. & Dar contexto a las direcciones observadas. & La asociación de ASN no reconstruye la ruta medida. \\
Docker Compose & Configuración y ejecución de contenedores. & Mantener servicios, red y almacenamiento definidos. & No elimina la variación propia de la red. \\
\bottomrule
\end{longtable}
\endgroup

En el flujo completo, primero obtuve los traceroutes con Scamper y RIPE Atlas. En paralelo conservé capturas para revisar qué se había enviado y recibido. Después normalicé las respuestas y añadí el contexto de ASN. Finalmente construí los grafos y los compartí en NDEx. Esta separación permite volver desde una figura publicada hasta las mediciones que la sustentan, sin confundir la visualización con la evidencia original.
